\uuid{Yn7O}
\titre{Calcul de dérivée}
\auteur{bodin ; hanani ; mzari}
\organisation{exo7}
\licence{CC BY-NC-SA 4.0}
\datecreate{2020-01-20}
\niveau{L1}
\module{Analyse}
\chapitre{Dérivabilité des fonctions réelles}
\sousChapitre{Autre}
\difficulte{1}
\theme{Dérivées, calcul de dérivée}
\exo7qcmid{lille-1.derivees.008}

\contenu{
\begin{question}
\qclassification{124.00}
\qtags{motcle=calcul de dérivée}
Soit $\displaystyle f(x)=x^2-\mathrm{e}^{x^2-1}$. Quelles sont les bonnes réponses ?
\begin{answers}
    \good{$f$ admet un minimum local en $0$.}
    \bad{$f$ admet un maximum local en $0$.}
    \bad{$f$ admet un point d'inflexion en $0$.}
    \bad{la tangente à $\mathcal{C}_f$ en $0$ est une droite verticale.}
\end{answers}
\begin{explanations}
On calcule $f'(x)=2x-2x\mathrm{e}^{x^2-1}$ et $f''(x)=2-2(1+2x^2)\mathrm{e}^{x^2-1}$. Ensuite, on vérifie que
$$f'(x)=0\mbox{ et }f''(0)=2-2\mathrm{e}^{-1}>0.$$
Donc $f$ admet un minimum local en $0$ et la tangente à $\mathcal{C}_f$ en $0$ est une droite horizontale.
\end{explanations}
\end{question}
}
